% Irodov Part III §3.2 (89--100), user PDF pp.112--114,
% printed pp.103--105. Independent Part-III numbering.
% Conventions: eps=relative permittivity; r_perp=distance from interface normal.
\electroconductproblem{89}{103}
真空半空间内有点电荷 $q$，距平面电介质界面为
$\ell$。介质占据另一半无限空间，相对介电常数为
$\varepsilon$。求：(a) 距电荷在界面上投影点
的径向距离为 $s$ 处，界面束缚电荷面密度；
讨论 $\ell\to0$；(b) 界面总束缚电荷。
\exsolution
取界面 $z=0$，真电荷位于 $(0,0,\ell)$。
介质边界的镜像法给真空区等效镜像电荷
$q'=q(1-\varepsilon)/(1+\varepsilon)$，
位于 $z=-\ell$。介质区的等效场相当于
$2q/(1+\varepsilon)$ 位于真电荷原位置。
介质侧的法向极化强度为
$\bm P=\varepsilon_0(\varepsilon-1)\bm E_{\rm diel}$。
在介质外法向为 $+\bm e_z$ 的边界，
介质内 $E_z=-[2q/(1+\varepsilon)]
\ell/[4\pi\varepsilon_0(s^2+\ell^2)^{3/2}]$，
于是
\[
\boxed{\sigma_b(s)=
-\frac{q\ell(\varepsilon-1)}
{2\pi(\varepsilon+1)(s^2+\ell^2)^{3/2}}}.
\]
固定 $s>0$ 时 $\ell\to0$ 有 $\sigma_b\to0$；
但在 $s=0$ 处呈集中奇性，
所以不能将总电荷误断为零。
积分平面：
\[
Q_b=\int_0^\infty2\pi s\,\sigma_b(s)\,\mathrm ds
=-\frac{q(\varepsilon-1)}{\varepsilon+1}
\ell\int_0^\infty\frac{s\,\mathrm ds}{(s^2+\ell^2)^{3/2}}
=\boxed{-q\frac{\varepsilon-1}{\varepsilon+1}}.
\]
\excheck 作为 $\varepsilon\to\infty$ 的导体极限，
$Q_b\to-q$，符合接地导体镜像解。

\electroconductproblem{90}{103}
利用上一题结果，求电介质界面束缚电荷对
位于界面上方 $\ell$ 处点电荷 $q$ 的作用力。
\exsolution
真空一侧的介质界面作用，可由位于镜像位置的
$q'=q(1-\varepsilon)/(1+\varepsilon)$
来表示，镜像距真电荷 $2\ell$。
因此作用力沿界面法线：
\[
F_z=\frac1{4\pi\varepsilon_0}
\frac{q q'}{(2\ell)^2}
=-\frac{q^2(\varepsilon-1)}
{16\pi\varepsilon_0\ell^2(\varepsilon+1)}.
\]
其模
\[
\boxed{F=\frac{q^2(\varepsilon-1)}
{16\pi\varepsilon_0\ell^2(\varepsilon+1)}}.
\]
当 $\varepsilon>1$，力指向介质。
\excheck $\varepsilon=1$ 时介质与真空不可区分，
力为零；$\varepsilon\to\infty$ 恢复导体镜像力。

\electroconductproblem{91}{104}
点电荷 $q$ 位于真空和均匀无限电介质
（相对介电常数 $\varepsilon$）的平面分界线上。
求两侧距电荷 $r$ 处的
$\bm D,\bm E$ 模及电势 $\varphi(r)$。
\exsolution
由界面法向位移与切向电场边界条件，
平面两侧的径向电场可取
$E(r)=A/r^2$，两侧取同一系数。
以电荷为球心的高斯球在真空、介质
各有一半球面积 $2\pi r^2$：
\[
q=2\pi r^2(\varepsilon_0E+
\varepsilon_0\varepsilon E)
=2\pi\varepsilon_0(1+\varepsilon)Ar^0.
\]
因此
\[
\boxed{E_{\rm vac}=E_{\rm diel}
=\frac{|q|}{2\pi\varepsilon_0(1+\varepsilon)r^2}}.
\]
位移矢量模为
\[
\boxed{D_{\rm vac}=\frac{|q|}
{2\pi(1+\varepsilon)r^2}},\qquad
\boxed{D_{\rm diel}=\frac{\varepsilon|q|}
{2\pi(1+\varepsilon)r^2}}.
\]
取无穷远电势零，径向积分得
\[
\boxed{\varphi(r)=
\frac{q}{2\pi\varepsilon_0(1+\varepsilon)r}}.
\]
\excheck 电势与 $\bm E$ 在两侧连续，
而 $\bm D$ 的模因介电常数不同而不同。

\electroconductproblem{92}{104}
在均匀电介质半空间内，距与真空分界面
$\ell$ 处有一个小导体球，带电 $q$。
介质相对介电常数 $\varepsilon$，
球半径远小于 $\ell$。
求平面束缚电荷密度的分布及 $\ell\to0$ 的行为。
\exsolution
在点电荷近似下，真电荷位于介质侧
$z=-\ell$。介质侧镜像电荷
$q'=q(\varepsilon-1)/(\varepsilon+1)$
位于 $z=+\ell$。
介质中的场为
$\bm E=(4\pi\varepsilon_0\varepsilon)^{-1}$
乘两点电荷叠加。
在界面 $z=0^-$，
两电荷的法向场为
\[
E_z=\frac{(q-q')\ell}
{4\pi\varepsilon_0\varepsilon(s^2+\ell^2)^{3/2}}
=\frac{q\ell}
{2\pi\varepsilon_0\varepsilon(\varepsilon+1)
(s^2+\ell^2)^{3/2}}.
\]
介质表面外法向为 $+\bm e_z$，所以
\[
\boxed{\sigma_b(s)=
\varepsilon_0(\varepsilon-1)E_z
=\frac{q\ell(\varepsilon-1)}
{2\pi\varepsilon(\varepsilon+1)
(s^2+\ell^2)^{3/2}}}.
\]
\excheck 对固定 $s>0$，$\ell\to0$ 时密度趋零；
当 $\ell$ 过小接近球半径时，点电荷近似本身失效。

\electroconductproblem{93}{104}
均匀电介质占据导体平面一侧半空间，
相对介电常数 $\varepsilon$。
距接地导体平面 $\ell$ 处有小金属球，
带电 $q$，球半径可忽略。
求导体边界处的介质束缚面电荷密度。
\exsolution
将导体平面视为 $z=0$，
真电荷位于 $z=\ell$，
镜像 $-q$ 位于 $z=-\ell$。
介质内部的法向场
$E_z(0^+)=-q\ell/
[2\pi\varepsilon_0\varepsilon(s^2+\ell^2)^{3/2}]$。
介质的外法向在底部是 $-\bm e_z$，
故
\[
\sigma_b=\bm P\cdot(-\bm e_z)
=-\varepsilon_0(\varepsilon-1)E_z(0^+).
\]
于是
\[
\boxed{\sigma_b(s)=
\frac{q\ell(\varepsilon-1)}
{2\pi\varepsilon(s^2+\ell^2)^{3/2}}}.
\]
\excheck 对正 $q$，靠近导体的介质束缚面
带正电，但导体感应\emph{自由电荷}为负电；
二者不能混为一谈。

\electroconductproblem{94}{104}
图3.15：两平行导体板短接，板距 $d$，
板间放一块沿板法向均匀固定极化的介质板，
厚度 $h<d$，极化强度为 $\bm P$。
求介质内部与剩余气隙中的 $\bm E,\bm D$。
\exsolution
取极化方向为正法向，用标量 $P$ 表示。
两个区域都没有自由体电荷，
法向位移恒为 $D$。
在介质板内部
$D=\varepsilon_0E_{\rm diel}+P$，
气隙内
$D=\varepsilon_0E_{\rm gap}$。
由于两极板短接，总电压为零：
\[
(d-h)E_{\rm gap}+hE_{\rm diel}=0.
\]
代入 $E_{\rm gap}=D/\varepsilon_0$、
$E_{\rm diel}=(D-P)/\varepsilon_0$ 得
$dD-hP=0$，因此
\[
\boxed{\bm D=\frac hd\bm P
\quad\text{（两区域）}},
\]
\[
\boxed{\bm E_{\rm gap}=\frac{h}{\varepsilon_0d}\bm P},
\qquad
\boxed{\bm E_{\rm diel}=
-\frac{1-h/d}{\varepsilon_0}\bm P}.
\]
\excheck $h=d$ 时介质完全填满，
气隙消失、内部电场为零；
$h\to0$ 时电介质作用趋零。

\electroconductproblem{95}{104}
无限长圆柱体电介质沿径向具有固定极化
$\bm P=\alpha\bm r_\perp$，
$\alpha>0$。求内部束缚体电荷密度。
\exsolution
束缚体电荷密度定义为
$\rho_b=-\nabla\cdot\bm P$。
在柱坐标里，
$P_r=\alpha r$，
\[
\nabla\cdot\bm P=
\frac1r\frac{\partial(rP_r)}{\partial r}
=\frac1r\frac{\partial(\alpha r^2)}{\partial r}
=2\alpha.
\]
所以
\[
\boxed{\rho_b=-2\alpha}.
\]
\excheck $\alpha$ 的单位为
$\mathrm{C/m^3}$，密度与半径无关；
若圆柱半径为 $R$，
侧表面还带束缚面电荷
$\sigma_b=\alpha R$。

\electroconductproblem{96}{104}
均匀固定极化的介质球，极化为 $\bm P$。
(a) 利用正负电荷的微小错位证明球内电场；
(b) 证明球外是中心点偶极子的电场，
并求其偶极矩。
\exsolution
极化均匀，束缚体密度
$\rho_b=-\nabla\cdot\bm P=0$，
球表面密度 $\sigma_b=P\cos\theta$。
其可看作密度为 $+\rho,-\rho$ 的两均匀带电球
相对微小偏移 $\bm\delta$，满足
$\rho\bm\delta=\bm P$。
在两球共同内部，场为
\[
\bm E=
\frac{\rho}{3\varepsilon_0}\bm r
-\frac{\rho}{3\varepsilon_0}(\bm r+\bm\delta)
=\boxed{-\frac{\bm P}{3\varepsilon_0}}.
\]
球外总电荷零，最低非零多极矩为偶极矩
\[
\boxed{\bm p=V\bm P=
\frac{4\pi R^3}{3}\bm P}.
\]
其球外电势可直接由
$\sigma_b=P\cos\theta$ 的球面积分验证为
\[
\boxed{\varphi(\bm r)=
\frac{\bm p\cdot\bm r}
{4\pi\varepsilon_0r^3}},\qquad r>R.
\]
\excheck 球内均匀退极化场指向
$-\bm P$；球外场不等于同一均匀场。

\electroconductproblem{97}{104}
无限均匀固定极化的电介质内部挖去
一个球形空腔。介质极化强度为 $\bm P$，
远离空腔的实际电场为 $\bm E_\infty$。
求空腔内电场 $\bm E_{\rm cav}$。
\exsolution
线性叠加：用一个均匀极化为 $-\bm P$ 的球
叠加到原均匀极化介质，
刚好抵消球内的极化，即挖去空腔。
根据第96题，一颗均匀极化为
$\bm P_{\rm sph}$ 的球在内部产生
$-\bm P_{\rm sph}/(3\varepsilon_0)$。
因此空腔附加场
$\Delta\bm E=-(-\bm P)/(3\varepsilon_0)
=+\bm P/(3\varepsilon_0)$，
所以
\[
\boxed{\bm E_{\rm cav}=
\bm E_\infty+\frac{\bm P}{3\varepsilon_0}}.
\]
\excheck 注意本题的远场 $\bm E_\infty$
是\emph{原介质中的实际宏观场}。
若原书简答出现相反符号，可能采用不同的
场变量约定；本式以叠加法的极化定义为准。

\electroconductproblem{98}{104}
均匀相对介电常数为 $\varepsilon$ 的
电介质球置于外加均匀场 $\bm E_0$，
球在作用下均匀极化。
利用第96题求球内电场与极化强度。
\exsolution
球内总电场为施加场与自身退极化场之和：
\[
\bm E=\bm E_0-\frac{\bm P}{3\varepsilon_0}.
\]
线性介质满足
$\bm P=\varepsilon_0(\varepsilon-1)\bm E$，
代入：
\[
\bm E\left(1+\frac{\varepsilon-1}{3}\right)
=\bm E_0.
\]
故
\[
\boxed{\bm E=\frac3{\varepsilon+2}\bm E_0},
\quad
\boxed{\bm P=
\frac{3\varepsilon_0(\varepsilon-1)}
{\varepsilon+2}\bm E_0}.
\]
\excheck $\varepsilon=1$ 时极化为零，
球内场等于外加场；
$\varepsilon\to\infty$ 时
球内场趋零。

\electroconductproblem{99}{105}
无限长圆柱形介质固定均匀极化，
极化 $\bm P$ 垂直于圆柱轴线。
求圆柱内部电场。
\exsolution
取圆柱轴沿 $z$、极化沿 $+x$，
侧表面束缚电荷密度
$\sigma_b=P\cos\phi$，内部无束缚体电荷。
二维边值问题的内外电势可以取
\[
\varphi_{\rm in}=Ar\cos\phi,\qquad
\varphi_{\rm out}=B\frac{\cos\phi}{r},
\]
其中表面电势连续给 $B=AR^2$。
在 $r=R$，法向场跃变
$E_{{\rm out},n}-E_{{\rm in},n}=
\sigma_b/\varepsilon_0$。
由
$E_{{\rm in},n}=-A\cos\phi$、
$E_{{\rm out},n}=+A\cos\phi$，
得 $2A=P/\varepsilon_0$。
因此
\[
\boxed{\bm E_{\rm in}=
-\frac{\bm P}{2\varepsilon_0}}.
\]
\excheck 与球体 $1/3$ 系数不同，
圆柱横向退极化因子是 $1/2$。

\electroconductproblem{100}{105}
无限长均匀介质圆柱置于均匀外场
$\bm E_0$ 内，圆柱轴线与 $\bm E_0$ 垂直。
圆柱相对介电常数为 $\varepsilon$。
利用上一题结果求柱内场与极化强度。
\exsolution
根据第99题，均匀极化产生
$-\bm P/(2\varepsilon_0)$ 退极化场，
因此
\[
\bm E=\bm E_0-\frac{\bm P}{2\varepsilon_0}.
\]
材料本构关系
$\bm P=\varepsilon_0(\varepsilon-1)\bm E$，
代入：
\[
\left(1+\frac{\varepsilon-1}{2}\right)\bm E
=\bm E_0,
\]
从而
\[
\boxed{\bm E=
\frac2{\varepsilon+1}\bm E_0},
\qquad
\boxed{\bm P=
\frac{2\varepsilon_0(\varepsilon-1)}
{\varepsilon+1}\bm E_0}.
\]
\excheck 当介质常数趋于1时，
内部场还原为外加场；
$\varepsilon\to\infty$ 时内部电场趋零。
